\section{Edge collection and the supplied-graph benchmark}

The characterization in \cref{obj:thm:observed-record-precision-frontier} gives calibrated edge retention a prospective interpretation: the known retention probability \(q\) is a collection choice whose precision consequences depend on population size and the common degree bound. Under the independent audit design in \cref{obj:ass:audit,obj:ass:design-independence}, increasing retention reduces the contribution of edge collection to squared risk. The comparison between \(d/(nq)\) and \(d^2/n\) places the transition at retention of order \(1/d\). At or above this order, the minimax risk has the same order as the full-retention scale \(\min\{1,d^2/n\}\). Below the transition, the collection term determines the uncapped risk scale.

Population size determines how this threshold translates into a requirement for consistent estimation. For the admissible sequences in \cref{obj:thm:observed-record-precision-frontier}, uniform consistency is equivalent to
\[
d_n=o(\sqrt n),
\qquad
q_n>0\ \text{eventually},
\qquad
\frac{nq_n}{d_n}\longrightarrow\infty.
\]
Thus a collection plan can use a vanishing retention probability provided it exceeds \(d_n/n\) by a diverging factor and the degree condition is satisfied. Retention of order \(1/d_n\) or greater suffices to match the full-retention precision order. When \(d_n^2/n\to0\), consistent collection plans can have retention below this transition. For a fixed degree bound, for example, consistency requires \(nq_n\to\infty\), while retention bounded away from zero attains the full-retention order.

The upper-risk envelope in \cref{obj:def:upper-envelope} supplies a complementary planning calculation. Its assignment component is \(5(d+1)^2/n\), and its audit component is \(4d(1-q)/(nq)\). The latter vanishes at full retention. Accordingly, the envelope makes the benefit of additional collection explicit, while the matching minimax characterization determines the attainable order over all measurable uses of the observed record.

To isolate the degree component, consider an experiment that supplies the true interference graph in addition to the original observations. The supplied-graph minimax squared risk \leanref{sym:RG}{\ensuremath{R^G_{n,d}(q)}} in \cref{obj:def:supplied-risk} retains the schedule class, causal target, assignment law, and estimator-randomization convention of the original decision problem.

\begin{definitionv}{def:supplied-risk}[Supplied-graph risk $R^G_{n,d}(q)$]
For fixed admissible \(n,d,q\), define the supplied-graph minimax squared risk by
\[
R^G_{n,d}(q)=
\inf_{T^G}\sup_{\theta\in\mathcal M_{n,d}}
\mathbb E_{P_{\theta,q}}
\left[
\bigl(T^G(O,G(\theta),\xi)-\tau(\theta)\bigr)^2
\right].
\]
Here \(\mathcal M_{n,d}\) is the schedule class, \(P_{\theta,q}\) is the fixed-schedule experiment law, \(O\) is the original observed record, \(\tau(\theta)\) is the population all-treated-versus-all-control contrast, and \(\xi\) is an independent uniform seed on \([0,1]\). The additional observation \(G(\theta)\) is the true off-diagonal graph component of the fixed schedule \(\theta\).

The infimum ranges over all Borel-measurable functions
\[
T^G:
\operatorname{range}(O)\times
2^{\{(j,i)\in V^2:j\ne i\}}\times[0,1]
\longrightarrow\mathbb R,
\]
where \(V\) is the fixed finite population. The finite graph coordinate carries the discrete sigma-algebra. The original record carries the Borel sigma-algebra inherited from its finite graph and assignment coordinates and its real outcome coordinates. The estimator may use the known \(n,d,q\). Expectations of the nonnegative loss take values in \([0,+\infty]\), and the estimator class includes every such measurable rule.
\end{definitionv}

Supplying the graph makes every true source-to-recipient arrow available to the estimator; the outcome coefficients and population contrast remain quantities to be learned from the randomized outcomes. The quadratic benchmark of \citet[Theorem \texttt{thm:degree-frontier}; non-archival research note]{CausalSmithKnownGraph2026} is translated to the present additive conventions in \cref{obj:thm:supported-boundaries}. After adjoining the compulsory diagonal coordinate as in \cref{obj:lem:published-model}, the supplied-graph score \(T_G\) is exactly additive, half-Bernoulli SNIPE: its own-treatment sign is the diagonal summand and its true-source signs are the off-diagonal summands. The present work adds inverse-inclusion correction for sampled arrows and analyzes that rule jointly with the unrestricted complete-record converse. Under the same schedule class in \cref{obj:def:model} and assignment law in \cref{obj:ass:assignment}, the supplied-graph lower bound and score calculation in \cref{obj:lem:supplied-block-record-lower,obj:lem:score-audit} give
\[
R^G_{n,d}(q)\asymp \min\left\{1,\frac{d^2}{n}\right\}
\]
uniformly over admissible \(n,d,q\), with universal comparison constants. The upper bound uses the supplied-graph score in \cref{obj:def:oracle-score}, clipped to the target range, together with the zero rule when appropriate. The lower bound identifies the same degree dependence even with the true graph supplied. This comparison therefore attributes the degree component to estimation under the fixed assignment law and schedule restrictions, while the additional retention dependence measures the precision cost of sampled graph information.

For a platform workflow based on sampled interaction records, this interpretation ties collection to explicit sampling and influence conditions. The experiment in \cref{obj:synth_3} requires independent arrow retention at a known common probability, independent half-Bernoulli treatment assignment, and observation of every assignment and realized outcome. Its influence conditions are the additive response model, both degree bounds, and the coefficient-mass restriction in \cref{obj:def:model}. The approximate-network analysis of \citet[Assumptions 1--3; Theorems 1--2; Lemma 2]{JiangDengWangWang2024ApproximateNetworks} instead gives variance \(O\{d_{\mathcal G}^2/[n\pi(1-\pi)]\}\) for its pseudoinverse estimator under bounded outcomes and first- and second-order influence restrictions, with an omitted-influence condition controlling relative bias. Here, the independent retention experiment supports inverse-inclusion correction, a complete-record converse, and exact joint consistency conditions, making population size, degree, and collection probability inputs to an unrestricted minimax planning result.
